\section*{Analytical Appendices: Financial Reasoning in the Field}

\subsection*{Appendix A: Reconcile the Example Through Journal Entries}

The 3-month prepaid-course case in Part V can be reconstructed with a small set of accounting entries, using debit and credit language conceptually:
\begin{enumerate}
  \item \textbf{January 1 (Capital contribution):} Debit Cash ₹5,00,000; Credit Share Capital ₹5,00,000.
  \item \textbf{January 1 (Advance tuition collected):} Debit Cash ₹3,00,000; Credit Contract Liability (Unearned Tuition) ₹3,00,000.
  \item \textbf{January 1 (Equipment purchased):} Debit Equipment Asset ₹60,000; Credit Cash ₹60,000.
  \item \textbf{January 15 to 31 (Marketing and delivery paid):} Debit Marketing Expense ₹1,20,000; Debit Teaching and Delivery Expense ₹60,000; Credit Cash ₹1,80,000.
  \item \textbf{January 31 (Revenue earned for Month 1):} Debit Contract Liability ₹1,00,000; Credit Course Revenue ₹1,00,000.
  \item \textbf{January 31 (Depreciation recognized):} Debit Depreciation Expense ₹5,000; Credit Accumulated Depreciation (Contra-Asset) ₹5,000.
\end{enumerate}

These entries are not tax advice or a full chart of accounts. They show how a founder can record a commercial event without confusing it with another event. The \textbf{advance} is not immediately revenue; the \textbf{asset purchase} is not automatically the same-period expense; the \textbf{equity contribution} is not a sale; and \textbf{depreciation} can reduce profit without consuming cash again that month. Proper bookkeeping combines these distinctions with contract-specific accounting, tax treatments, and controls.

\subsection*{Appendix B: A Receivable Case That Turns Profit into a Cash Crisis}

Suppose a technical services company completes an engagement and recognizes ₹10 lakh of revenue. It has ₹6 lakh of expenses appropriately recognized in the same period, leaving ₹4 lakh of simplified accounting profit before tax and financing. But the client has not yet paid. The business still needs to pay ₹5 lakh of its expenses in cash immediately and the remaining ₹1 lakh next month. It may report profit and face acute liquidity pressure at the same time.

If it began with only ₹2 lakh of cash and has no other collections, it cannot pay the ₹5 lakh immediate obligation without borrowing, delaying suppliers, obtaining customer advances, or changing payment terms. The problem is not necessarily poor unit contribution; it is financing the operating cycle. A founder who mistakes recognized revenue for bank cash may accept a large contract that endangers the company.

To prevent this, attach collection milestones, expected dates, approval conditions, invoice requirements, and downside timing to every institutional deal. Forecast cash \textbf{before} signing, using the actual employer, college, or enterprise payment process rather than the founder's preferred schedule.

\subsection*{Appendix C: A Simple LTV Sensitivity Table}

Suppose a recurring product has contribution of ₹500 per retained customer-month, received at month-end, with 1 percent monthly discounting ($d = 0.01$) and a constant monthly survival probability $r$. In the stylized model from Section 25, expected lifetime contribution before CAC is:
\begin{equation}
LTV_{\text{before CAC}} = \frac{₹500}{1.01 - r}.
\end{equation}
Table~\ref{tab:ltv-sensitivity} displays the extreme sensitivity of expected lifetime contribution to monthly survival rates.

\begin{table}[htbp]
\centering
\small
\begin{tabularx}{\columnwidth}{Xr}
\toprule
\textbf{Monthly Survival Probability ($r$)} & \textbf{Lifetime Contribution Before CAC} \\
\midrule
0.70 & ₹1,613 \\
0.80 & ₹2,381 \\
0.90 & ₹4,545 \\
0.95 & ₹8,333 \\
\bottomrule
\end{tabularx}
\caption{Sensitivity of expected customer lifetime contribution to monthly survival rate.}
\label{tab:ltv-sensitivity}
\end{table}

The arithmetic shows why small changes in assumed retention can produce enormous swings in LTV. It does \textbf{not} prove those survival rates exist or that retention will remain constant. With a tiny early sample, forecasting 95 percent monthly survival for years is often unjustified. Build cohorts, report observed horizons and uncertainty, and refuse to finance acquisition with a retention estimate chosen because it makes the investment attractive.

\subsection*{Appendix D: The Funnel-to-Unit-Economics Bridge}

Imagine a venture spends ₹60,000 on an acquisition test and reaches 600 relevant people. 90 respond, 30 qualify, 12 receive a priced offer, and 6 pay. The total test CAC is \textbf{₹10,000 per paid customer} before any acquisition labor not included in the ₹60,000. If price is ₹25,000 and direct variable delivery cost is ₹7,000, contribution before CAC is ₹18,000 and after the test CAC is ₹8,000 per buyer. Across 6 buyers, this is ₹48,000 before cohort fixed costs and overhead.

Now imagine the same channel receives more reach but qualified conversion halves while the marketing expenditure stays constant. CAC doubles to ₹20,000. The unit contribution after acquisition becomes \textbf{negative ₹2,000} at the same price and delivery cost. The underlying educational product did not change; the business model deteriorated because the distribution economics changed.

A founder must follow the funnel to paid customers and then through delivered outcomes. A ``cheap lead'' is attractive only if enough leads become suitable, paying, supported customers at sustainable cost.
